\section{Discussion and limitations}\label{sec:disc}
\textbf{Answer to the question.} With labels, unlabeled data and predictor fixed, and the predictor output as the only information, an interval can be shorter than PPI++ only to the extent that $\mathbb E[Y\mid f]$ is not affine, by at most $1-\sqrt{1-sN/(n+N)}$ in width (Cor.~\ref{cor:gap}), and only if $n$ is large enough to learn the curve. On the seven data sets that ship with \texttt{ppi-python} the unmodified predictors are close to calibrated ($s\le2.5\%$), so the available gain is $\le1.2\%$ in width and cross-fitted recalibration did not realise it at $n\ge100$ (it cost 0.5--2\% at $n=500$ and up to 19\% at $n=50$). For predictors whose scores are monotone but not on the scale of the label the gain is real (2--27\% here) and is predicted accurately by the bound; there the improvement over PPI ($\lambda=1$) is much larger still. A hard 0/1 predictor leaves no room. Recalibration also removes the dependence on the scale of the scores, which PPI/PPI++ do not have. A practical rule that follows from the theory: plot the labeled $(f,Y)$ calibration curve, estimate $s$ and the bound by cross-validation, and recalibrate only if $(N/m)\,s$ is of the order of the estimation cost ($\approx$ a few percent at $n\approx500$); and use the unclipped PPI++ weight, which costs nothing and avoided losses of 2--6\% on two data sets.

\textbf{Limitations.}
(1) \emph{Target.} Only the population mean was studied; the closed-form gap uses the mean's influence function. For quantiles or regression coefficients the efficient estimators involve other conditional functions, and our numerical findings do not transfer without new checks.
(2) \emph{Information set.} The bound is for estimators that use only the scalar predictor output. Features $X$ could help in principle; we only measured population-level gains with gradient boosting on the three census sets (negligible), not interval-level behaviour, and learning from $X$ is closer to training a new predictor than to the fixed-predictor constraint of the task.
(3) \emph{Labeling mechanism.} Labeled points are an i.i.d.\ random subset (MCAR). Under shift (plankton) the recalibrated intervals were much less robust than PPI; we did not study covariate-shift corrections.
(4) \emph{Asymptotics and small $n$.} Intervals are Wald-type; Prop.~\ref{prop:valid} is asymptotic and not proved for the CV-selected learner (\textbf{auto}). Coverage at $n\le100$ is below nominal for all methods on heavy-tailed outcomes (Census income: 82--92\%) and can be 1--4 points low even on benign data; we do not claim to fix that, and ``shorter'' claims are only meaningful where coverage is nominal. A $t$-quantile did not repair PPI++ at $n=50$.
(5) \emph{Resampled populations.} Real-data cells resample the empirical distribution of finite example sets; for Forest (1596 records) and for $n$ comparable to $P$, repeated records make the effective problem easier for nonparametric fits. The distortions in Sec.~\ref{sec:exp-distort} are constructed, so we have no estimate of how often real deployed predictors are non-affinely calibrated; the two data sets with the largest $s$ (Gene expression 2.5\%, Census income 0.5\%) are those whose predictors are regressions on continuous outcomes, which is suggestive but not a sample of predictors.
(6) \emph{Novelty.} The efficiency bound and the cross-fitted augmented estimator are classical results in semiparametric/semi-supervised statistics; our contributions are the closed-form comparison with PPI and PPI++, the validity statement without rates, and the empirical map of when the gap is nonzero. Related prior work on the same estimand may exist that we did not find.
(7) \emph{Scope of evidence.} 2000 (synthetic) or 1000 (real) replicates per cell give coverage standard errors of about 0.5--0.7 points, so differences below one point in coverage are not resolved; width ratios are paired and much more precise.
\textbf{Broader impact.} The work is methodological, uses public data and has no foreseeable direct negative societal impact; the main risk is overconfidence from intervals used outside their assumptions (labeling mechanism, shift), which the failure cases document.
